% answers_short.tex —— 1991 年·三南卷 · 简短答案速查
% 由 tools/make_short_answers.py 自 answers.tex 的 \gksolhead/\gkanswer 提取，
% 并与 1991/1991_数学真题_三南_答案要点.md 逐题比对后生成，勿手改（改源 MD 后重跑）。

\begin{gknotice}
\textbf{使用说明}：本节只给最终答案便于快速核对，完整推导见解析版。答案为编写者独立推导，\textbf{非官方标准答案}。
\end{gknotice}

\gkpart{选择题}
\begin{center}\begin{tabular}{ccccccccccccccccccc}
\toprule
题号 & 1 & 2 & 3 & 4 & 5 & 6 & 7 & 8 & 9 & 10 & 11 & 12 & 13 & 14 & 15 & 16 & 17 & 18 \\
\midrule
答案 & B & A & C & B & C & B & D & D & C & A & D & B & B & D & B & A & C & D \\
\bottomrule
\end{tabular}\end{center}

\gkpart{填空题}
\begin{itemize}[leftmargin=2.6em]
  \item[\textbf{19.}] $\dfrac{\sqrt7}4$
  \item[\textbf{20.}] $1$
  \item[\textbf{21.}] $\dfrac{4r}3$
  \item[\textbf{22.}] $0$
  \item[\textbf{23.}] $\dfrac V3-V_1$
  \item[\textbf{24.}] $0$
\end{itemize}

\gkpart{解答题}
\begin{itemize}[leftmargin=2.6em]
  \item[\textbf{25.}] $\cos\beta=\dfrac{3\sqrt{10}}{10}$
  \item[\textbf{26.}] $\left[-5,-1+\dfrac{\sqrt{14}}2\right]$
  \item[\textbf{27.}] $AB_1\perp A_1M$，证明见解析。
  \item[\textbf{28.}] $a_n=\dfrac{n+1}2$，$b_n=6\left(\dfrac13\right)^{n-1}$，$n\ge1$。
  \item[\textbf{29.}] 以焦点中点为原点、焦轴为 $x$ 轴的正交坐标系下，$x^2-\dfrac{y^2}3=1$。原题未预设坐标系。
  \item[\textbf{30.}] (1) $f(x)$ 在 $\mathbb R$ 上严格递增；(2) 所有整数 $n\ge3$ 均满足 $f(n)>\dfrac n{n+1}$
\end{itemize}
